%==============================================================================
\section{Bias and Variance --- Độ lệch và phương sai mô hình}
\label{sec:bias-variance}
%==============================================================================

\begin{dongco}[title={Hai nguồn sai số chính}]
Bias và variance mô tả nguồn lỗi trong dự đoán của mô hình ML.
\textbf{Bias}: lỗi do \textbf{đơn giản hoá quá mức} quan hệ giữa feature và đầu ra.
\textbf{Variance}: lỗi do \textbf{nhạy quá} với dao động nhỏ trong dữ liệu huấn luyện.
Mục tiêu: giảm cả hai để dự đoán tốt trên dữ liệu chưa thấy.
\end{dongco}

\begin{ynghia}[title={Ba thành phần lỗi}]
Bias, variance và \textbf{irreducible error} (nhiễu không giảm được).
Có \textbf{đánh đổi} bias--variance: giảm bias thường tăng variance và ngược lại.
\end{ynghia}

\subsection{Bias là gì?}

\begin{dinhnghia}[title={Bias}]
Bias được tính là chênh lệch giữa \textbf{dự đoán trung bình} và \textbf{giá trị thực}.
Trong ML, bias (sai số hệ thống) xảy ra khi mô hình \textbf{giả định sai} về dữ liệu.
\end{dinhnghia}

\begin{ynghia}[title={High bias vs low bias}]
\begin{itemize}
  \item \textbf{High bias}: mô hình quá đơn giản, không khớp train lẫn test $\Rightarrow$ \textbf{underfitting}.
  \item \textbf{Low bias}: bắt được pattern ẩn; nếu quá phức tạp dễ dẫn tới variance cao và overfitting.
\end{itemize}
\end{ynghia}

\begin{vidu}[title={Minh hoạ bias cao / thấp}]
\tsimg{07_bias_variance/img01.jpg}
\end{vidu}

\begin{phuongphap}[title={Ví dụ mô hình}]
Hồi quy tuyến tính trên dữ liệu phi tuyến $\Rightarrow$ bias cao.
Linear/logistic regression thường bias cao hơn; cây quyết định, kNN, SVM thường bias thấp hơn (linh hoạt hơn).
\end{phuongphap}

\subsection{Variance là gì?}

\begin{dinhnghia}[title={Variance}]
Trong thống kê, variance đo độ trải của tập số quanh mean.
Trong ML, variance là \textbf{biến thiên dự đoán} khi huấn luyện trên các tập dữ liệu hơi khác nhau.
\end{dinhnghia}

\begin{ynghia}[title={High variance vs low variance}]
\begin{itemize}
  \item \textbf{High variance}: khớp cả nhiễu trên train, kém trên test $\Rightarrow$ \textbf{overfitting}.
  \item \textbf{Low variance}: ít thay đổi giữa train/test; mô hình quá đơn giản có thể kèm bias cao.
\end{itemize}
\end{ynghia}

\begin{vidu}[title={Minh hoạ variance cao / thấp}]
\tsimg{07_bias_variance/img02.jpg}
\end{vidu}

\subsection{Đánh đổi Bias--Variance}

\begin{ynghia}[title={Cân bằng}]
Tăng độ phức tạp mô hình $\Rightarrow$ bias giảm, variance tăng.
Giảm độ phức tạp $\Rightarrow$ bias tăng, variance giảm.
Cần điểm tối ưu sao cho \textbf{tổng lỗi dự đoán} nhỏ nhất.
\end{ynghia}

\begin{vidu}[title={Đồ thị đánh đổi Bias--Variance}]
\tsimg{07_bias_variance/img03.jpg}

Trục $X$: độ phức tạp mô hình; trục $Y$: lỗi dự đoán.
Tổng lỗi = bias + variance (+ irreducible); vùng tối ưu là chỗ cân bằng hai thành phần.
\end{vidu}

\begin{ynghia}[title={Biểu thức toán học}]
\[
\text{Error} = \text{bias}^2 + \text{variance} + \text{irreducible error}.
\]
Giảm lỗi dự đoán bằng cách chọn độ phức tạp cân bằng bias và variance.
\end{ynghia}

\begin{phuongphap}[title={Kỹ thuật cân bằng}]
\textbf{Giảm bias cao}: mô hình phức tạp hơn; thêm feature; giảm regularization.

\textbf{Giảm variance cao}: regularization; đơn giản hoá mô hình; thêm dữ liệu; cross-validation.
\end{phuongphap}

\subsection{Ví dụ Python}

\begin{vidu}[title={High bias --- Linear Regression}]
\begin{lstlisting}
import numpy as np
from sklearn.datasets import fetch_california_housing
from sklearn.model_selection import train_test_split
from sklearn.linear_model import LinearRegression
from sklearn.metrics import mean_squared_error

data = fetch_california_housing()
X, y = data.data, data.target
X_train, X_test, y_train, y_test = train_test_split(
    X, y, test_size=0.2, random_state=42)

lr = LinearRegression()
lr.fit(X_train, y_train)

train_mse = mean_squared_error(y_train, lr.predict(X_train))
test_mse = mean_squared_error(y_test, lr.predict(X_test))
print("Training MSE:", train_mse)
print("Testing MSE:", test_mse)
\end{lstlisting}

\textbf{Output} (mức MSE phụ thuộc dữ liệu; train và test đều cao, gợi ý bias đáng kể so với mô hình phức tạp hơn).
\end{vidu}

\begin{vidu}[title={Low bias, high variance --- Polynomial Regression}]
\begin{lstlisting}
from sklearn.preprocessing import PolynomialFeatures

poly = PolynomialFeatures(degree=2)
X_train_poly = poly.fit_transform(X_train)
X_test_poly = poly.transform(X_test)

pr = LinearRegression()
pr.fit(X_train_poly, y_train)

train_mse = mean_squared_error(y_train, pr.predict(X_train_poly))
test_mse = mean_squared_error(y_test, pr.predict(X_test_poly))
print("Training MSE:", train_mse)
print("Testing MSE:", test_mse)
\end{lstlisting}

Train MSE thường \textbf{thấp hơn nhiều} test MSE $\Rightarrow$ overfitting / variance cao.
\end{vidu}

\begin{vidu}[title={Giảm variance --- Ridge Regression}]
\begin{lstlisting}
from sklearn.linear_model import Ridge

ridge = Ridge(alpha=1)
ridge.fit(X_train_poly, y_train)

train_mse = mean_squared_error(y_train, ridge.predict(X_train_poly))
test_mse = mean_squared_error(y_test, ridge.predict(X_test_poly))
print("Training MSE:", train_mse)
print("Testing MSE:", test_mse)
\end{lstlisting}

Test MSE thường cải thiện so với đa thức không regularization --- variance giảm, bias tăng nhẹ.
\end{vidu}

\begin{vidu}[title={Tối ưu $\alpha$ --- GridSearchCV}]
\begin{lstlisting}
from sklearn.model_selection import GridSearchCV

param_grid = {'alpha': np.logspace(-3, 3, 7)}
ridge_cv = GridSearchCV(Ridge(), param_grid, cv=5)
ridge_cv.fit(X_train_poly, y_train)

train_mse = mean_squared_error(y_train, ridge_cv.predict(X_train_poly))
test_mse = mean_squared_error(y_test, ridge_cv.predict(X_test_poly))
print("Training MSE:", train_mse)
print("Testing MSE:", test_mse)
\end{lstlisting}

Chọn $\alpha$ qua CV để cân bằng bias--variance tốt hơn.
\end{vidu}
